\chapter{Field-model options for LIONS}\label{ch:options}

\section{Case H: which option removes the excess, at what cost}\label{sec:caseH}
\textbf{Question.}
Which opt-in additions to LIONS remove the G1 and G2 excess errors, and at what cost per particle-step?
The bilinear field and the Boris step stay the default, and the Fortran bit-identity of the default path is kept.

\textbf{Design} (registered before running).
Four field models on the same node data (F0 bilinear, F1 bicubic Hermite, F2 cubic and F3 quintic splines;
\cref{ch:interpolation}) and two integrators: Boris, and RK4 with a fixed step.
Three step counts (720, 1440, 2880 per turn), and three maps: 3-sector and 10-sector FROM\_ORBIT at \thirddeg{} and
\SI{1}{\degree}, and the PSI Ring.
Decision target, fixed before the run: $|\Delta\Qr|$, $|\Delta\Qz|$ and $|\det M - 1|$ all at most $10^{-4}$ on every orbit
or energy of a map.
The validation came first.
F1 reproduces node values and derivatives to \HFOneNodes, and RK4 shows fourth order from 360 steps per turn.

\textbf{Result} (\cref{fig:H}, \cref{tab:Hdecision,tab:Hcost}).
\begin{itemize}
\item F0 never meets the target, with either integrator.
RK4 does not fix it: on the 10-sector and PSI maps its worst $|\det M - 1|$ at 1440 steps is \HFzeroRkDetMin{} to
\HFzeroRkDetMax, \HBorisOverRkMin{} to \HBorisOverRkMax{} times below Boris on all maps.
Both decrease with the step count, slowly: fixed-step integration across the gradient jumps limits every integrator.
\item F1 with Boris meets the target wherever any option does, at the lowest cost per turn, on the FROM\_ORBIT maps:
720 steps on the 3-sector map, 2880 on the 10-sector map at \thirddeg{}.
\item No option meets it on the 10-sector map at \SI{1}{\degree}: too few nodes per harmonic period
(\cref{ch:interpolation}).
\item On the PSI map the reference is RF-Track on F3, which favours F3; only F3 meets the target against it.
Against FIXPO, the external reference, F1 with Boris at 1440 steps is the closest model: median $|\Delta\Qr|$
\HFixpoHermQrMed{} and $|\Delta\Qz|$ \HFixpoHermQzMed, against \HFixpoBilQrMed{} and \HFixpoBilQzMed{} for F0,
\HFixpoCubQrMed{} and \HFixpoCubQzMed{} for F2, \HFixpoQuinQrMed{} and \HFixpoQuinQzMed{} for F3.
FIXPO's own interpolation is unknown; that F1 shares its use of the file derivatives is a hypothesis.
\end{itemize}

\begin{figure}[t]
\centering
\includegraphics[width=\textwidth]{h_matrix}
\caption{Case H: worst error over orbits or energies (largest of $|\Delta\Qr|$, $|\Delta\Qz|$, $|\det M-1|$) against
steps per turn, for each field model and integrator.
Dotted line: decision target $10^{-4}$.
PSI: reference RF-Track on F3; the 386 to \SI{430}{MeV} band is excluded.}
\label{fig:H}
\end{figure}

\begin{table}[h]
\centering\small
\caption{Case H: smallest step count per turn (of 720, 1440, 2880) meeting the decision target; -- means none.
B: Boris, RK: RK4.}
\label{tab:Hdecision}
\input{gen/tab_h_decision.tex}
\end{table}

\begin{table}[h]
\centering\small
\caption{Case H: cost per particle-step in the numpy harness, $10^5$ particles, PSI map, Apple M3 CPU (ratio to F0 with
Boris).
LIONS itself costs \HLionsCost{} times the harness F0, so the harness cost is representative.}
\label{tab:Hcost}
\input{gen/tab_h_cost.tex}
\end{table}

\section{F1 in LIONS}
F1 is in LIONS as the deck option \texttt{interpolation=hermite} (API: \texttt{Deck.field\_interpolation}).
LIONS derives the Hermite slopes from its own map: $\partial_rB$ and $\partial_\theta B$ from $\partial_xB$, $\partial_yB$, and the
cross derivative as in F1.
So it needs no new input file.
The Fortran code ignores the option token, which is checked by a test.

Validation against the harness: one turn agrees within \HPhaseTwoOneTurn{} relative.
Closed orbits, turn times and $\Qz$ agree to rounding.
$\Qr$ and $\det M - 1$ differ by up to \HPhaseTwoJacMax.
The registered tolerance was $10^{-9}$, so the validation failed as written.
The residual is attributed to the finite-difference Jacobian acting on the remaining rounding noise (not tested), and
it is \HTargetRatio{} or more times below the decision target.
The cost is \HPhaseTwoCostRatio{} times the default on the PSI map at $10^5$ particles.

The first implementation took the azimuth from $\arccos$ (\cref{sec:azimuth}).
Its Jacobian differed from the harness by an amount that grew as $1/h$ with the probe step $h$.
That is the signature of non-smooth noise in the one-turn map, not of a field difference.
Taking $\theta$ from $\operatorname{atan2}$ removed most of it.

\section{The \texttt{azimuth=atan2} option}
The same change is available for the default bilinear path, as \texttt{azimuth=atan2}.
\Cref{tab:Hazimuth} compares the radial Jacobian of LIONS with that of the numpy bilinear tracker (which uses
$\operatorname{atan2}$).

\begin{table}[h]
\centering\small
\caption{Maximum $|\Delta M_{rr}|$ between LIONS bilinear and the numpy bilinear tracker, for three Jacobian steps.}
\label{tab:Hazimuth}
\input{gen/tab_h_azimuth.tex}
\end{table}

The registered criterion, a tenfold reduction at $h = \SI{1e-6}{m}$, was not met on either map.
On the 10-sector map at $h = \SI{1e-7}{m}$ the noise even grew.
The working hypothesis, not tested, is that the bilinear one-turn map has a second noise source: it is only piecewise
smooth, so probes that straddle a cell edge mix two slopes.
On the C$^1$ Hermite field the same azimuth change worked, which fits this hypothesis.

\section{What to use}
\begin{itemize}
\item For tune or Jacobian work on measured maps: \texttt{interpolation=hermite}.
Its effect on autograd gradients has not been measured: the PyTorch port is bilinear only.
It lowers $|\det M-1|$ by \HDetRatioMin{} to \HDetRatioMax{} times at 1440 steps, for \HCostHerm{} times the cost of
the default; on the PSI map it still leaves \HFoneDetPsi{} at 1440 steps and \HFoneDetPsiTwoEightEighty{} at 2880 steps, while the
tunes do not meet the target against the F3 reference at any step count.
\item Check the node density per harmonic period of the map (\cref{ch:interpolation}) before trusting any interpolation.
\item The default stays for reproducing Fortran results.
\end{itemize}
